Для примера входных данных
\begin{verbatim}
4
2 1
3 3
4 4
5 2
\end{verbatim}

возможен следующий вывод:

\begin{verbatim}
6
2 0
2 3
5 3
5 2
4 2
4 4
\end{verbatim}

\begin{center}
\includegraphics[scale=0.5]{1.png}
\end{center}

Обратите внимание, что данный пример не фигурирует среди тестов задачи.

\textbf{Visualizer}

Среди приложений к этой задаче имеется скрипт, позволяющий визуализировать
входные и выходные файлы.

Чтобы визуализировать входной файл, используйте следующую команду:

\texttt{python vis.py [input file]}

Также вы можете визуализировать ваш ответ для некоторого теста, используя
следующую команду. По техническим причинам визуализатор показывает только
\texttt{первые $1000$ отрезков} из выходного файла.


\texttt{python vis.py [input file] --solution [output file]}

Пример:

\texttt{python vis.py examples/00.in --solution examples/00.out}
